%!TEX root = ../../main.tex
%% ===============================================================
%% 9.1 기준 성능
%% 파일: chapters/09_results/01_baseline.tex
%% 상위: chapters/09_results/chapter.tex
%% 정의 라벨: sec:results-baseline
%% 참조 라벨: tab:baseline
%% 포함 객체: tables/tab_baseline
%% 인용: -
%% 상태: 초안 (docs/OUTLINE.md 참고)
%% ===============================================================

\section{기준 성능}
\label{sec:results-baseline}

%% ---------------------------------------------------------------
%% table 객체: tab:baseline — 입력 표현별 시험 AUC
%% 파일: tables/tab_baseline.tex   사용처: chapters/09_results/01_baseline.tex
%% 수치는 config/numbers.tex 매크로만 사용 (재실행 시 numbers.tex 만 갱신)
%% ---------------------------------------------------------------
\begin{table}[htbp]
  \centering
  \caption{입력 표현별 시험 AUC (패치 15.16, 세 시드 평균)}
  \label{tab:baseline}
  \small
  \begin{tabular}{@{}llc@{}}
    \toprule
    모델 & 입력 표현 & 시험 AUC \\
    \midrule
    \textbf{LightGBM} & 공학적 테이블 요약 (약 2{,}980 차원) & \textbf{\aucLGBM{}} \\
    MLP & 동일 테이블 요약 (정합 입력) & \aucMLP{} \\
    BiGRU & 거시 시퀀스 (95 차원 $\times$ 6) & \aucBiGRU{} \\
    계층적 융합 & 전역 + 그래프 + 사건 & \aucLayered{} \\
    Transformer & 거시 시퀀스 & \aucTransformer{} \\
    Cross-Attn & 사건 토큰 $\times$ 플레이어 그래프 & \aucXAttn{} \\
    ST-GNN & 시공간 플레이어 그래프 & \aucSTGNN{} \\
    GraphSAGE & 플레이어 그래프 & \aucGraphSAGE{} \\
    \bottomrule
  \end{tabular}
\end{table}


표 \ref{tab:baseline}에서 두 가지 구조적 사실이 드러난다. 첫째, 공학적 테이블 요약 위의 LightGBM이 $\aucLGBM{}$로 가장 높고, 정합 입력 MLP가 $\aucMLP{}$으로 그 다음이다. 둘째, 신경망 기반의 시퀀스·그래프·사건 표현은 $.569$--$\aucBiGRU{}$의 좁은 띠에 모여 있으며, 세 스트림을 모두 결합한 계층적 융합도 단일 BiGRU와 같은 $\aucBiGRU{}$에 머문다. 즉 표현을 늘리는 것이 신호를 더 드러내지 못했다.
